\section{On-Policy Distillation: Basics}
\label{sec:on-policy-distillation-basics}

On-policy distillation (OPD) computes supervision on trajectories sampled from
the current student policy $\pi_\theta$. Given a prompt $x\sim\Dcal_x$, the
student samples a response
\begin{align*}
    y &\sim \pi_\theta(\cdot\mid x),
    \qquad T(y) \coloneqq |y| \text{ (the length of the response)},\\
    y &= (y_1,\ldots,y_{T(y)}),
\end{align*}
where the end-of-sequence token is included in the response. Both the student
and the fixed teacher $\pi_{\mathrm{Teacher}}$ are evaluated on each
student-generated prefix $y_{<t}$. Their next-token distributions are, for all $v\in\Vcal$ (the vocabulary),
\begin{align*}
    p_t(v)
    \coloneqq \pi_\theta(v\mid x,y_{<t}),
    \quad\text{and}\quad
    q_t(v)
    \coloneqq \pi_{\mathrm{Teacher}}(v\mid x,y_{<t}).
\end{align*}

A standard formulation minimizes the sequence-level reverse KL on the prompt
distribution. Let $\Ycal(x)$ denote the set of complete responses to $x$, and
define the joint distribution $\mu_\theta$ by sampling
$x\sim\Dcal_x$ first and then $y\sim\pi_\theta(\cdot\mid x)$. In the discrete
case, this means
$\mu_\theta(x,y)=\Dcal_x(x)\pi_\theta(y\mid x)$. Expanding the KL before
combining the expectations gives
\begin{align}
    \Lcal_{\mathrm{OPD}}(\theta)
    &=
    \E_{x\sim\Dcal_x}
    \left[
        D_{\mathrm{KL}}\!\left(
            \pi_\theta(\cdot\mid x)
            \,\middle\|\,
            \pi_{\mathrm{Teacher}}(\cdot\mid x)
        \right)
    \right]
    \nonumber\\
    &=
    \E_{x\sim\Dcal_x}
    \left[
        \sum_{y\in\Ycal(x)}
        \pi_\theta(y\mid x)
        \log\frac{\pi_\theta(y\mid x)}
                 {\pi_{\mathrm{Teacher}}(y\mid x)}
    \right]
    \nonumber\\
    &=
    \E_{x\sim\Dcal_x}
    \left[
        \E_{y\sim\pi_\theta(\cdot\mid x)}
        \left[
            \log\frac{\pi_\theta(y\mid x)}
                     {\pi_{\mathrm{Teacher}}(y\mid x)}
        \right]
    \right]
    \nonumber\\
    &=
    \E_{(x,y)\sim\mu_\theta}
    \left[
        \log\frac{\pi_\theta(y\mid x)}
                 {\pi_{\mathrm{Teacher}}(y\mid x)}
    \right].
    \label{eq:opd.sequence-kl}
\end{align}
The last line is the law of total expectation. It does not make $x$ and $y$
independent: their joint law is $\mu_\theta$, under which the response $y$ is
conditioned on the sampled prompt $x$.

\begin{highlight}
    \eqref{eq:opd.sequence-kl} can be equivalently rewritten as the token-level form loss function
    \begin{align*}
        \Lcal_{\mathrm{OPD}}(\theta)
        =
        \E_{(x,y)\sim\mu_\theta}
        \left[
            \sum_{t=1}^{T(y)}
            D_{\mathrm{KL}}\!\left(
                \pi_\theta(\cdot\mid x,y_{<t})
                \,\middle\|\,
                \pi_{\mathrm{Teacher}}(\cdot\mid x,y_{<t})
            \right)
        \right].
    \end{align*}
\end{highlight}

To obtain the token-level form, fix a prompt $x$ and abbreviate the student
and teacher sequence probabilities by $P_\theta(y\mid x)$ and
$Q_{\mathrm{Teacher}}(y\mid x)$. Autoregressive factorization gives
\begin{align}
    P_\theta(y\mid x)
    &=
    \prod_{t=1}^{T(y)}\pi_\theta(y_t\mid x,y_{<t}),
    &
    Q_{\mathrm{Teacher}}(y\mid x)
    &=
    \prod_{t=1}^{T(y)}
    \pi_{\mathrm{Teacher}}(y_t\mid x,y_{<t}).
    \label{eq:opd.factorization}
\end{align}
Substituting these products into the definition of KL and using
$\log\prod_t a_t=\sum_t\log a_t$ yields
\begin{align}
    &D_{\mathrm{KL}}\!\left(
        P_\theta(\cdot\mid x)
        \,\middle\|\,
        Q_{\mathrm{Teacher}}(\cdot\mid x)
    \right)
    \nonumber\\
    &\quad=
    \sum_y P_\theta(y\mid x)
    \log\frac{P_\theta(y\mid x)}{Q_{\mathrm{Teacher}}(y\mid x)}
    = 
    \E_{y\sim\pi_\theta(\cdot\mid x)}
    \left[
        \log
        \frac{
            \prod_{t=1}^{T(y)}\pi_\theta(y_t\mid x,y_{<t})
        }{
            \prod_{t=1}^{T(y)}
            \pi_{\mathrm{Teacher}}(y_t\mid x,y_{<t})
        }
    \right]
    \nonumber\\
    &\quad=
    \E_{y\sim\pi_\theta(\cdot\mid x)}
    \left[
        \sum_{t=1}^{T(y)}
        \log
        \frac{
            \pi_\theta(y_t\mid x,y_{<t})
        }{
            \pi_{\mathrm{Teacher}}(y_t\mid x,y_{<t})
        }
    \right].
    \label{eq:opd.log-ratio-sum}
\end{align}
For a fixed $x$ and a realized prefix $y_{<t}$, define the corresponding conditional reverse KL by
\begin{align}
    \mathrm{kl}_t(x,y_{<t})
    &\coloneqq
    \E_{v\sim\pi_\theta(\cdot\mid x,y_{<t})}
    \left[\log\frac{\pi_\theta(v\mid x,y_{<t})}{\pi_{\mathrm{Teacher}}(v\mid x,y_{<t})}\right]
    \nonumber\\
    &=
    D_{\mathrm{KL}}\!\left(
        \pi_\theta(\cdot\mid x,y_{<t})
        \,\middle\|\,
        \pi_{\mathrm{Teacher}}(\cdot\mid x,y_{<t})
    \right).
    \label{eq:opd.conditional-token-kl}
\end{align}

Assume that generation is capped at $T_{\max}$. For each $t$, the indicator
$\mathbf{1}\{t\leq T(y)\}$ is determined by $y_{<t}$ because it only records
whether an end-of-sequence token has already appeared. Therefore it can be
pulled through the conditional expectation over the next token. Using the
tower property term by term gives
\begin{align}
    &\E_{y\sim\pi_\theta(\cdot\mid x)}
    \left[
        \sum_{t=1}^{T(y)}
        \log\frac{\pi_\theta(y_t\mid x,y_{<t})}{\pi_{\mathrm{Teacher}}(y_t\mid x,y_{<t})}
    \right]
    \nonumber\\
    &\quad=
    \sum_{t=1}^{T_{\max}}
    \E_{y\sim\pi_\theta(\cdot\mid x)}
    \left[
        \mathbf{1}\{t\leq T(y)\}
        \log\frac{\pi_\theta(y_t\mid x,y_{<t})}{\pi_{\mathrm{Teacher}}(y_t\mid x,y_{<t})}
    \right]
    \nonumber\\
    &\quad=
    \sum_{t=1}^{T_{\max}}
    \E_{y\sim\pi_\theta(\cdot\mid x)}
    \left[
        \mathbf{1}\{t\leq T(y)\}
        \E_{v\sim\pi_\theta(\cdot\mid x,y_{<t})}
        \left[\log\frac{\pi_\theta(v\mid x,y_{<t})}{\pi_{\mathrm{Teacher}}(v\mid x,y_{<t})}\right]
    \right]
    \nonumber\\
    &\quad=
    \E_{y\sim\pi_\theta(\cdot\mid x)}
    \left[
        \sum_{t=1}^{T(y)}
        \mathrm{kl}_t(x,y_{<t})
    \right].
    \label{eq:opd.tower-decomposition}
\end{align}
The second equality follows by conditioning each outer expectation on the
realized prefix $y_{<t}$. Conditional on $x$ and $y_{<t}$, the actual next
token has distribution
$y_t\sim\pi_\theta(\cdot\mid x,y_{<t})$. Its conditional mean is therefore
$\E_{v\sim\pi_\theta(\cdot\mid x,y_{<t})} \left[\log\frac{\pi_\theta(v\mid x,y_{<t})}{\pi_{\mathrm{Teacher}}(v\mid x,y_{<t})}\right]$, where $v$ is only a dummy variable for the next token and the resulting quantity depends only on $y_{<t}$;
averaging it under the complete-response distribution $y\sim
\pi_\theta(\cdot\mid x)$ simply averages over the student-induced prefixes.
The unobserved suffix integrates to one.

Combining \eqref{eq:opd.log-ratio-sum},
\eqref{eq:opd.conditional-token-kl}, and
\eqref{eq:opd.tower-decomposition}, then restoring the prompt expectation,
yields
\begin{align}
    \Lcal_{\mathrm{OPD}}(\theta)
    =
    \E_{(x,y)\sim\mu_\theta}
    \left[
        \sum_{t=1}^{T(y)}
        D_{\mathrm{KL}}\!\left(
            \pi_\theta(\cdot\mid x,y_{<t})
            \,\middle\|\,
            \pi_{\mathrm{Teacher}}(\cdot\mid x,y_{<t})
        \right)
    \right].
    \label{eq:opd.token-kl}
\end{align}

\subsection{Implemnted OPD Objectives}
In practice, the actual loss computation varies with different implementations on how this exact per-token reverse KL is computed. We will see some examples below.

\begin{itemize}
    \item \textbf{Sampled-token OPD.}
    The lightest variant evaluates the log-ratio only at the token $y_t$
    sampled by the student. At a realized prefix $(x,y_{<t})$, define
    \begin{align}
        \ell_t^{\mathrm{sample}}(x,y_{\leq t})
        \coloneqq
        \log p_t(y_t)-\log q_t(y_t).
        \label{eq:opd.sampled-token-loss}
    \end{align}
    Its token-level objective is
    \begin{highlight}
        \begin{align}
            \Lcal_{\mathrm{OPD}}^{\mathrm{sample}}(\theta)
            =
            \E_{(x,y)\sim\mu_\theta}
            \left[
                \sum_{t=1}^{T(y)}
                \ell_t^{\mathrm{sample}}(x,y_{\leq t})
            \right].
            \label{eq:opd.sampled-token-objective}
        \end{align}
    \end{highlight}
    Since $y_t\sim p_t$ conditional on $(x,y_{<t})$,
    \begin{align}
        \E_{y_t\sim p_t}
        \left[
            \ell_t^{\mathrm{sample}}(x,y_{\leq t})
        \right]
        =
        \sum_{v\in\Vcal}p_t(v)\log\frac{p_t(v)}{q_t(v)}
        =
        D_{\mathrm{KL}}(p_t\|q_t).
        \label{eq:opd.sampled-token-unbiased}
    \end{align}
    Thus, the sampled-token loss is an unbiased one-sample estimator of the
    scalar token-level reverse KL at the sampled prefix.

    \item \textbf{Full-vocabulary OPD.}
    This variant evaluates the divergence over the entire vocabulary at every
    sampled prefix:
    \begin{highlight}
        \begin{align}
            \Lcal_{\mathrm{OPD}}^{\mathrm{full}}(\theta)
            =
            \E_{(x,y)\sim\mu_\theta}
            \left[
                \sum_{t=1}^{T(y)}
                D_{\mathrm{KL}}(p_t\|q_t)
            \right].
            \label{eq:opd.full-vocabulary-objective}
        \end{align}
    \end{highlight}
    It provides dense supervision but materializing both distributions costs
    $O(BT_{\max}\Vcal)$ ( where $|\Vcal|$ is the vocabulary size) memory for batch size $B$.

    \item \textbf{Student-top-$k$ OPD.}
    This intermediate variant restricts the divergence to the $k$ tokens with
    highest student probability:
    \begin{align*}
        S_t\coloneqq S_t(x,y_{<t})
        \coloneqq
        \operatorname{TopK}(p_t,k)
        \subseteq \Vcal.
    \end{align*}
    Both distributions are renormalized on this same support:
    \begin{align}
        \bar p_t^{(S_t)}(v)
        &\coloneqq
        \frac{p_t(v)\mathbf{1}\{v\in S_t\}}
             {\sum_{u\in S_t}p_t(u)},
        &
        \bar q_t^{(S_t)}(v)
        &\coloneqq
        \frac{q_t(v)\mathbf{1}\{v\in S_t\}}
             {\sum_{u\in S_t}q_t(u)}.
        \label{eq:opd.top-k-renormalization}
    \end{align}
    The resulting objective is
    \begin{highlight}
        \begin{align}
            \Lcal_{\mathrm{OPD}}^{\mathrm{top}\text{-}k}(\theta)
            =
            \E_{(x,y)\sim\mu_\theta}
            \left[
                \sum_{t=1}^{T(y)}
                D_{\mathrm{KL}}\!\left(
                    \bar p_t^{(S_t)}
                    \,\middle\|\,
                    \bar q_t^{(S_t)}
                \right)
            \right].
            \label{eq:opd.top-k-objective}
        \end{align}
    \end{highlight}

    \begin{remark}
        \item
        \begin{itemize}
        \item Probability mass outside $S_t$ is discarded, so this is generally not
        equal to the full-vocabulary reverse KL. It preserves multi-token
        supervision on the student's high-probability region at lower cost. In other words, the sampled-token OPD is unbiased yet suffers from high variance while the full-vocabulary OPD is unbiased and has lower variance, yet it requires more memory. Top-k OPD is a compromise between the two, it is biased but has lower variance than sampled-token OPD. It also loses the probability mass information outside $S_t$.

        \item \textbf{\eqref{eq:opd.top-k-objective} is ill-defined when $k=1$.} When $k=1$, this objective seems to degenerate. If $S_t = \{v^*_t\}$, then both restricted distributions become point masses after renormalization $\bar p_t(v^*_t) = 1, \bar q_t(v^*_t) = 1$, so the KL is identically zero. This would provide no distillation signal.
        
        \item    \textbf{When is top-$k$ OPD a good surrogate for full-vocabulary OPD?}
        Define the mass each distribution places on the student top-$k$ set,
        $P_{S_t}=\sum_{v\in S_t}p_t(v)$ and $Q_{S_t}=\sum_{v\in S_t}q_t(v)$.
        Write $S_t^c=\Vcal\setminus S_t$ for the complement, and assume
        $0<P_{S_t},Q_{S_t}<1$ so that both pieces below are proper distributions.
        Renormalize on the complement in the same way as
        \eqref{eq:opd.top-k-renormalization}:
        \begin{align}
            p_t^{S_t^c}(v)
            &\coloneqq
            \frac{p_t(v)\mathbf{1}\{v\in S_t^c\}}{1-P_{S_t}},
            &
            q_t^{S_t^c}(v)
            &\coloneqq
            \frac{q_t(v)\mathbf{1}\{v\in S_t^c\}}{1-Q_{S_t}}.
            \label{eq:opd.complement-renormalization}
        \end{align}
        Partition the reverse KL into the top-$k$ support and its complement:
        \begin{align}
            D_{\mathrm{KL}}(p_t\|q_t)
            &=
            \sum_{v\in\Vcal}p_t(v)\log\frac{p_t(v)}{q_t(v)}
            \nonumber\\
            &=
            \sum_{v\in S_t}p_t(v)\log\frac{p_t(v)}{q_t(v)}
            +
            \sum_{v\in S_t^c}p_t(v)\log\frac{p_t(v)}{q_t(v)}.
            \label{eq:opd.kl-split}
        \end{align}
        On $S_t$, the renormalization
        \eqref{eq:opd.top-k-renormalization} says
        $p_t(v)=P_{S_t}\bar p_t^{(S_t)}(v)$ and
        $q_t(v)=Q_{S_t}\bar q_t^{(S_t)}(v)$. Substituting and using
        $\sum_{v\in S_t}\bar p_t^{(S_t)}(v)=1$ separates the shape of the
        conditional distribution from the mass on the set:
        \begin{align}
            \sum_{v\in S_t}p_t(v)\log\frac{p_t(v)}{q_t(v)}
            &=
            \sum_{v\in S_t}
            P_{S_t}\bar p_t^{(S_t)}(v)
            \log\frac{P_{S_t}\bar p_t^{(S_t)}(v)}{Q_{S_t}\bar q_t^{(S_t)}(v)}
            \nonumber\\
            &=
            P_{S_t}
            \sum_{v\in S_t}
            \bar p_t^{(S_t)}(v)
            \left(
                \log\frac{\bar p_t^{(S_t)}(v)}{\bar q_t^{(S_t)}(v)}
                +
                \log\frac{P_{S_t}}{Q_{S_t}}
            \right)
            \nonumber\\
            &=
            P_{S_t}\,
            D_{\mathrm{KL}}\!\left(
                \bar p_t^{(S_t)}
                \,\middle\|\,
                \bar q_t^{(S_t)}
            \right)
            +
            P_{S_t}\log\frac{P_{S_t}}{Q_{S_t}}.
            \label{eq:opd.kl-on-support}
        \end{align}
        The same substitution on the complement,
        $p_t(v)=(1-P_{S_t})p_t^{S_t^c}(v)$ and
        $q_t(v)=(1-Q_{S_t})q_t^{S_t^c}(v)$, gives
        \begin{align}
            \sum_{v\in S_t^c}p_t(v)\log\frac{p_t(v)}{q_t(v)}
            &=
            (1-P_{S_t})\,
            D_{\mathrm{KL}}\!\left(
                p_t^{S_t^c}
                \,\middle\|\,
                q_t^{S_t^c}
            \right)
            +
            (1-P_{S_t})\log\frac{1-P_{S_t}}{1-Q_{S_t}}.
            \label{eq:opd.kl-on-complement}
        \end{align}
        The two scalar terms are the KL between Bernoulli random variables that
        only record whether the next token falls in $S_t$. If
        $Z\sim\mathrm{Bern}(P_{S_t})$ and $W\sim\mathrm{Bern}(Q_{S_t})$, then
        \begin{align}
            D_{\mathrm{KL}}(Z\|W)
            &=
            P_{S_t}\log\frac{P_{S_t}}{Q_{S_t}}
            +
            (1-P_{S_t})\log\frac{1-P_{S_t}}{1-Q_{S_t}}.
            \label{eq:opd.bernoulli-kl}
        \end{align}
        Adding \eqref{eq:opd.kl-on-support} and \eqref{eq:opd.kl-on-complement},
        and identifying \eqref{eq:opd.bernoulli-kl}, produces the decomposition
        \begin{align}
            D_{\mathrm{KL}}(p_t\|q_t)
            &=
            P_{S_t}\,
            D_{\mathrm{KL}}\!\left(
                \bar p_t^{(S_t)}
                \,\middle\|\,
                \bar q_t^{(S_t)}
            \right)
            +
            (1-P_{S_t})\,
            D_{\mathrm{KL}}\!\left(
                p_t^{S_t^c}
                \,\middle\|\,
                q_t^{S_t^c}
            \right)
            \nonumber\\
            &\quad+
            D_{\mathrm{KL}}\!\left(
                \mathrm{Bern}(P_{S_t})
                \,\middle\|\,
                \mathrm{Bern}(Q_{S_t})
            \right).
            \label{eq:opd.kl-decomposition}
        \end{align}
        Top-$k$ OPD keeps only the renormalized divergence on $S_t$, and it
        drops the weight $P_{S_t}$. The other two terms are the bias. The
        complement term vanishes as $P_{S_t}\to 1$, which is the regime in which
        the student's top-$k$ set already carries essentially all of its mass.
        The Bernoulli term is small when $P_{S_t}\approx Q_{S_t}$, i.e., when
        the teacher places a similar fraction of its mass on that same set. In
        that regime $P_{S_t}\approx 1$ as well, so dropping the weight changes
        the objective by little. Outside it, top-$k$ OPD ignores both how the
        leftover mass is shaped and how much total mass each model assigns to
        $S_t$.
        \end{itemize}
    \end{remark}
\end{itemize}

\subsection{Gradients of the OPD Objective}
\label{sec:opd.gradients}

The three variants above describe what is \emph{computed}; this section describes what
is \emph{differentiated}. We first derive the exact gradient of
\eqref{eq:opd.token-kl}, then show that standard OPD keeps only part of it, and finally
check, variant by variant, whether plain backpropagation produces that part.

Throughout, the teacher is fixed, so $\nabla_\theta q_t=0$. We use one identity
repeatedly: for any probability $\pi_\theta(z)>0$ that is differentiable in $\theta$,
the chain rule gives
\begin{align}
    \pi_\theta(z)\,\nabla_\theta\log\pi_\theta(z)
    =
    \pi_\theta(z)\,\frac{\nabla_\theta\pi_\theta(z)}{\pi_\theta(z)}
    =
    \nabla_\theta\pi_\theta(z).
    \label{eq:opd.log-derivative}
\end{align}
We also use the stop-gradient operator $\operatorname{sg}[\cdot]$ (\texttt{.detach()} in code), with
the following convention. $\operatorname{sg}$ appears only inside an expression that is still to
be differentiated, i.e., a loss, a surrogate, or an advantage that enters a loss, where it marks a
factor that is treated as a constant under $\nabla_\theta$. An expression that already contains
$\nabla_\theta$ never carries $\operatorname{sg}$: every factor outside $\nabla_\theta$ is simply a
value at the current $\theta$.

\paragraph{The exact gradient.}
In \eqref{eq:opd.token-kl}, $\theta$ enters in two places: inside each per-state KL loss
$\mathrm{kl}_t(x,y_{<t})$ defined in \eqref{eq:opd.conditional-token-kl}, through $\pi_\theta(\cdot\mid x,y_{<t})$, and through $\mu_\theta$, which determines which prefixes $y_{<t}$ are visited. The prompt
distribution $\Dcal_x$ does not depend on $\theta$, so fix $x$ and differentiate
\begin{align*}
    J_x(\theta)
    & \coloneqq
    \E_{y\sim\pi_\theta(\cdot\mid x)}
    \left[
        \sum_{t=1}^{T(y)}
    D_{\mathrm{KL}}\!\left(
        \pi_\theta(\cdot\mid x,y_{<t})
        \,\middle\|\,
        \pi_{\mathrm{Teacher}}(\cdot\mid x,y_{<t})\right)
    \right].
    \nonumber\\
    & =
    \E_{y\sim\pi_\theta(\cdot\mid x)}
    \left[
        \sum_{t=1}^{T(y)}\mathrm{kl}_t(x,y_{<t})
    \right].
\end{align*}
\begin{enumerate}
    \item \emph{Reduce each term to a sum over prefixes.} As in the derivation of
    \eqref{eq:opd.tower-decomposition}, cap generation at $T_{\max}$ and use the indicator
    $\mathbf{1}\{t\leq T(y)\}$, which is determined by $y_{<t}$. The $t$-th summand then
    depends on $y$ only through $y_{<t}$, so its expectation only needs the marginal
    probability of the prefix,
    $\pi_\theta(y_{<t}\mid x)=\prod_{s=1}^{t-1}p_s(y_s)$:
    \begin{align*}
        J_x(\theta)
        =
        \sum_{t=1}^{T_{\max}}
        \sum_{y_{<t}}
        \pi_\theta(y_{<t}\mid x)\,
        \mathbf{1}\{t\leq T(y)\}\,
        \mathrm{kl}_t(x,y_{<t}).
    \end{align*}
    \item \emph{Product rule.} Every sum is finite, so the gradient passes inside. Each
    summand is a product of two $\theta$-dependent factors, $\pi_\theta(y_{<t}\mid x)$ and
    $\mathrm{kl}_t(x,y_{<t})$; the indicator does not depend on $\theta$. Applying the
    product rule, then \eqref{eq:opd.log-derivative} to $\nabla_\theta\pi_\theta(y_{<t}\mid x)$,
    \begin{align*}
        \nabla_\theta J_x(\theta)
        &=
        \sum_{t=1}^{T_{\max}}
        \sum_{y_{<t}}
        \mathbf{1}\{t\leq T(y)\}
        \left[
            \pi_\theta(y_{<t}\mid x)\,\nabla_\theta\mathrm{kl}_t(x,y_{<t})
            +
            \mathrm{kl}_t(x,y_{<t})\,\nabla_\theta\pi_\theta(y_{<t}\mid x)
        \right]
        \\
        &=
        \sum_{t=1}^{T_{\max}}
        \sum_{y_{<t}}
        \pi_\theta(y_{<t}\mid x)\,
        \mathbf{1}\{t\leq T(y)\}
        \left[
            \nabla_\theta\mathrm{kl}_t(x,y_{<t})
            +
            \mathrm{kl}_t(x,y_{<t})\,\nabla_\theta\log\pi_\theta(y_{<t}\mid x)
        \right].
    \end{align*}
    Undoing step 1 (the bracket still depends on $y$ only through $y_{<t}$) turns this
    back into an expectation over complete responses:
    \begin{align*}
        \nabla_\theta J_x(\theta)
        =
        \E_{y\sim\pi_\theta(\cdot\mid x)}
        \left[
            \sum_{t=1}^{T(y)}\nabla_\theta\mathrm{kl}_t(x,y_{<t})
        \right]
        +
        \E_{y\sim\pi_\theta(\cdot\mid x)}
        \left[
            \sum_{t=1}^{T(y)}\mathrm{kl}_t(x,y_{<t})\,\nabla_\theta\log\pi_\theta(y_{<t}\mid x)
        \right].
    \end{align*}
    \item \emph{Expand the prefix score.} Taking the log of
    $\pi_\theta(y_{<t}\mid x)=\prod_{s=1}^{t-1}p_s(y_s)$ gives
    $\nabla_\theta\log\pi_\theta(y_{<t}\mid x)=\sum_{s=1}^{t-1}\nabla_\theta\log p_s(y_s)$.
    The second expectation therefore contains the double sum
    \begin{align*}
        \sum_{t=1}^{T(y)}\sum_{s=1}^{t-1}
        \mathrm{kl}_t(x,y_{<t})\,\nabla_\theta\log p_s(y_s).
    \end{align*}
    \item \emph{Swap the order of summation.} The double sum runs over all pairs
    $1\leq s<t\leq T(y)$. Summing over $s$ on the outside instead,
    \begin{align*}
        \sum_{t=1}^{T(y)}\sum_{s=1}^{t-1}
        \mathrm{kl}_t(x,y_{<t})\,\nabla_\theta\log p_s(y_s)
        =
        \sum_{s=1}^{T(y)}
        \nabla_\theta\log p_s(y_s)
        \sum_{t=s+1}^{T(y)}\mathrm{kl}_t(x,y_{<t}).
    \end{align*}
    Renaming $s\to t$ and $t\to t'$ and restoring the expectation over $x\sim\Dcal_x$ gives
    \begin{align}
        \nabla_\theta\Lcal_{\mathrm{OPD}}(\theta)
        =
        \underbrace{
        \E_{(x,y)\sim\mu_\theta}
        \left[
            \sum_{t=1}^{T(y)}\nabla_\theta \mathrm{kl}_t(x,y_{<t})
        \right]
        }_{\text{(A): current state}}
        +
        \underbrace{
        \E_{(x,y)\sim\mu_\theta}
        \left[
            \sum_{t=1}^{T(y)}\nabla_\theta\log p_t(y_t)
            \sum_{t'=t+1}^{T(y)}\mathrm{kl}_{t'}(x,y_{<t'})
        \right]
        }_{\text{(B): future states}}.
        \label{eq:opd.exact-gradient}
    \end{align}
\end{enumerate}
Term (A) changes $\pi_\theta(\cdot\mid x,y_{<t})$ to reduce the KL at the state already visited. Term (B) is a
REINFORCE term whose ``reward'' for token $y_t$ is the total KL at all later positions:
$y_t$ determines which future states are visited, so it is credited with their KL. Note
that $\mathrm{kl}_t$ itself does not appear in the reward of $y_t$, because
$\mathrm{kl}_t(x,y_{<t})$ does not depend on $y_t$.

\paragraph{What standard OPD keeps.}
Standard OPD drops term (B): each visited state $(x,y_{<t})$ is treated as a constant
and only term (A) is kept. In reinforcement-learning terms, this corresponds to a
discount factor of zero. The three variants differ only in how they compute (or
estimate) $\nabla_\theta\mathrm{kl}_t(x,y_{<t})$, or its top-$k$ analogue, at a fixed state.
Below, write $\delta_t(v)\coloneqq\log p_t(v)-\log q_t(v)$, so that
$\mathrm{kl}_t(x,y_{<t})=\sum_{v\in\Vcal}p_t(v)\,\delta_t(v)$.

\paragraph{Full-vocabulary and top-$k$ OPD: plain backpropagation is correct.}
The full-vocabulary variant computes $\mathrm{kl}_t(x,y_{<t})=\sum_{v\in\Vcal}p_t(v)\,\delta_t(v)$
exactly. The top-$k$ variant computes
$D_{\mathrm{KL}}(\bar p_t^{(S_t)}\|\bar q_t^{(S_t)})$; the renormalized log-probabilities
are obtained by subtracting a log-sum-exp,
$\log p_t(i)-\log\sum_{j\in S_t}p_t(j)=\log\bar p_t^{(S_t)}(i)$ by
\eqref{eq:opd.top-k-renormalization}, and likewise for the teacher. The index set $S_t$
is chosen from detached log-probabilities, so it is held fixed; this is correct
because $S_t$ is piecewise constant in $\theta$. In both variants the value is a
deterministic, differentiable function of $\theta$ at a fixed state, with no sampling
inside, so autograd returns its exact gradient. For the top-$k$ variant this is the
exact gradient of the top-$k$ objective, which is itself a biased stand-in for
$\mathrm{kl}_t$ (see the remark after \eqref{eq:opd.top-k-objective}).

\paragraph{Sampled-token OPD: plain backpropagation is wrong.}
Here the computed value $\ell_t^{\mathrm{sample}}=\delta_t(y_t)$ depends on the token
$y_t\sim p_t$, which is drawn from the distribution being trained. Its expectation is
the target, $\E_{y_t\sim p_t}[\delta_t(y_t)]=\mathrm{kl}_t(x,y_{<t})$ by
\eqref{eq:opd.sampled-token-unbiased}, but autograd cannot differentiate through the
sampling step. The argument is a one-step version of the derivation above.
\begin{enumerate}
    \item \emph{What autograd computes.} Since $\nabla_\theta q_t=0$, backpropagating
    $\ell_t^{\mathrm{sample}}$ gives $\nabla_\theta\delta_t(y_t)=\nabla_\theta\log p_t(y_t)$.
    The teacher does not appear.
    \item \emph{Its expectation is zero.} Because $\Vcal$ is finite, the expectation is
    a finite sum, so
    \begin{align}
        \E_{y_t\sim p_t}\!\left[\nabla_\theta\log p_t(y_t)\right]
        =
        \sum_{v\in\Vcal}p_t(v)\,\nabla_\theta\log p_t(v)
        \overset{\eqref{eq:opd.log-derivative}}{=}
        \sum_{v\in\Vcal}\nabla_\theta p_t(v)
        =
        \nabla_\theta\sum_{v\in\Vcal}p_t(v)
        =
        \nabla_\theta 1
        =
        0 
        \label{eq:opd.sampled-token-gradient-expectation}
    \end{align}
    So the update is zero-mean noise that carries no information from the teacher.
    \item \emph{The true gradient.} Applying the product rule inside the finite sum,
    \begin{align*}
        \nabla_\theta\mathrm{kl}_t(x,y_{<t})
        &=
        \nabla_\theta\sum_{v\in\Vcal}p_t(v)\,\delta_t(v)
        =
        \sum_{v\in\Vcal}\nabla_\theta p_t(v)\,\delta_t(v)
        +
        \sum_{v\in\Vcal}p_t(v)\,\nabla_\theta\delta_t(v)
        \\
        &\overset{\eqref{eq:opd.log-derivative}}{=}
        \underbrace{\sum_{v\in\Vcal}p_t(v)\,\delta_t(v)\,\nabla_\theta\log p_t(v)}_{\E_{y_t\sim p_t}[\delta_t(y_t)\nabla_\theta\log p_t(y_t)]}
        +
        \underbrace{\sum_{v\in\Vcal}p_t(v)\,\nabla_\theta\log p_t(v)}_{=\,0\text{ by \eqref{eq:opd.sampled-token-gradient-expectation}}}.
    \end{align*}
    All of the signal is in the first sum, which comes from differentiating the sampling
    distribution. We record this as
    \begin{align}
        \nabla_\theta\mathrm{kl}_t(x,y_{<t})
        =
        \sum_{v\in\Vcal}p_t(v)\,\delta_t(v)\,\nabla_\theta\log p_t(v).
        \label{eq:opd.full-kl-grad}
    \end{align}
    \item \emph{The surrogate.} The pseudocode below therefore returns the loss as 
    $\operatorname{sg}[\delta_t(y_t)]\log p_t(y_t)$ instead of $\delta_t(y_t)$. Its value
    has no meaning, but its gradient is $\delta_t(y_t)\nabla_\theta\log p_t(y_t)$, whose
    expectation over $y_t\sim p_t$ is the first sum in step 3, i.e.,
    $\nabla_\theta\mathrm{kl}_t(x,y_{<t})$. This is the score-function (REINFORCE)
    estimator, and it is the same construction as $A_t^{\mathrm{OPD}}$ in PG-OPD implemented in VeRL (disscussed below).
\end{enumerate}

\subsection{Pseudocode}

The pseudocode below implements the three variants. In \texttt{full} and \texttt{topk},
the returned value is the per-state divergence itself; in \texttt{sampled}, it is the
surrogate from Section~\ref{sec:opd.gradients}.

\begin{codeblock}[language=Python]
def opd_token_loss(
    student_logits,
    teacher_logits,
    mode="full",
    sampled_tokens=None,
    top_k=None,
):
    """Return one OPD loss per batch element and response position."""
    s_logp = student_logits.float().log_softmax(dim=-1)
    with torch.no_grad():
        t_logp = teacher_logits.float().log_softmax(dim=-1)

    if mode == "sampled":
        if sampled_tokens is None:
            raise ValueError("sampled_tokens is required")
        token_idx = sampled_tokens.unsqueeze(dim=-1)
        s_token = s_logp.gather(dim=-1, index=token_idx).squeeze(dim=-1)
        t_token = t_logp.gather(dim=-1, index=token_idx).squeeze(dim=-1)
        # Not the loss value itself; see the gradient discussion above.
        return (s_token - t_token).detach() * s_token

    if mode == "full":
        return (s_logp.exp() * (s_logp - t_logp)).sum(dim=-1)

    if mode == "topk":
        if top_k is None:
            raise ValueError("top_k is required")

        # Student-selected support S_t.
        idx = s_logp.detach().topk(top_k, dim=-1).indices
        s_top = s_logp.gather(dim=-1, index=idx)
        t_top = t_logp.gather(dim=-1, index=idx)

        # Renormalize student and teacher on the same support.
        s_bar = s_top - s_top.logsumexp(dim=-1, keepdim=True)
        t_bar = t_top - t_top.logsumexp(dim=-1, keepdim=True)
        return (s_bar.exp() * (s_bar - t_bar)).sum(dim=-1)

    raise ValueError(f"unknown OPD mode: {mode}")


teacher = load_teacher()  # A model stronger than the student
student = init_student()
optimizer = AdamW(student.parameters(), lr=1e-6)

for step in range(total_steps):
    prompt = sample(prompts)

    # The student generates the states on which distillation is performed.
    with torch.no_grad():
        response_tokens = student.generate(prompt, max_tokens=max_len)

    trajectory = concatenate(prompt, response_tokens)
    s_all_logits = student(trajectory).logits
    with torch.no_grad():
        t_all_logits = teacher(trajectory).logits

    # Select next-token logits aligned with y_1, ..., y_{T(y)} only.
    s_logits = response_position_logits(s_all_logits, prompt)
    t_logits = response_position_logits(t_all_logits, prompt)

    # Three alternatives; choose one for a training run.
    sampled_loss = opd_token_loss(s_logits, t_logits, mode="sampled", sampled_tokens=response_tokens)
    # full_loss = opd_token_loss(s_logits, t_logits, mode="full")
    # topk_loss = opd_token_loss(s_logits, t_logits, mode="topk", top_k=64)

    # Exclude padding after each response's EOS token.
    active = response_activity_mask(response_tokens)
    loss = (full_loss * active).sum() / active.sum()
    optimizer.zero_grad()
    loss.backward()
    optimizer.step()
\end{codeblock}


\section{OPD Implementation Discussions}

\subsection{Original VeRL}
VeRL has two paths for you to use OPD:
\begin{itemize}
    \item \textbf{PG-OPD}: The negative distillation loss is used as the token-level advantage estimation in the PPO loss. Specifically, for the sampled student token $y_t$, construct the detached token-level signal:
    $$ 
    A_t^{\mathrm{OPD}} = -\operatorname{sg}\left[ \log p_t(y_t) - \log q_t(y_t)\right], 
    $$
    where $\operatorname{sg}$ is the stop-gradient operator, $p_t(v)
    \coloneqq \pi_\theta(v\mid x,y_{<t}), \text{ and }
    q_t(v)\coloneqq \pi_{\mathrm{Teacher}}(v\mid x,y_{<t})$.
    Then feed it into the existing \href{https://github.com/verl-project/verl/blob/adc7eefa16dad75c5f7b878823d5a76eac90c7b3/verl/trainer/distillation/losses.py#L280-L288}{PPO policy-loss calculation}.
    
    \item \textbf{GKD-OPD}: The distillation loss is the target. However, the VeRL implements the \citet{agarwal2024policy} so the forward-KL is used, i.e., when using the top-k to do the KL estimation, the teacher-top-k support is used.
\end{itemize}

\subsubsection{OPD as a PPO Update}
\label{sec:opd.as-ppo}

The VeRL actor is built for PPO. Given responses sampled from the rollout policy
$\pi_{\theta_{\mathrm{old}}}$ and one advantage $\hat A_t$ per response position, it minimizes the
clipped objective, aggregated over response positions according to \texttt{loss\_agg\_mode},
\begin{align}
    \max\!\Big\{
        -\hat A_t\,w_t(\theta),\;
        -\hat A_t\operatorname{clip}\!\big(w_t(\theta),1-\epsilon_{\mathrm{low}},1+\epsilon_{\mathrm{high}}\big)
    \Big\},
    \qquad
    w_t(\theta)
    \coloneqq
    \frac{p_t(y_t)}{p_t^{\mathrm{old}}(y_t)},
    \label{eq:opd.ppo-clip}
\end{align}
where $p_t^{\mathrm{old}}(v)\coloneqq\pi_{\theta_{\mathrm{old}}}(v\mid x,y_{<t})$. PG-OPD reuses this
objective unchanged and only supplies the advantage: in the
\href{https://github.com/verl-project/verl/blob/adc7eefa16dad75c5f7b878823d5a76eac90c7b3/verl/trainer/distillation/losses.py#L280-L288}{code},
\icode{advantages=-distillation_losses.detach()}, where \texttt{distillation\_losses} is the
per-token log-ratio (the \texttt{k1} estimator). In the notation of
Section~\ref{sec:opd.gradients},
\begin{align*}
    \hat A_t=A_t^{\mathrm{OPD}}=-\operatorname{sg}\!\left[\delta_t(y_t)\right].
\end{align*}

\paragraph{The update loop.}
Each PPO iteration of the \href{https://github.com/verl-project/verl/blob/adc7eefa16dad75c5f7b878823d5a76eac90c7b3/verl/trainer/ppo/ray_trainer.py#L1560}{trainer}
samples responses with the current parameters $\theta_{\mathrm{old}}$, computes
\texttt{old\_log\_probs} once, i.e., $\log p_t^{\mathrm{old}}(y_t)$, and then calls the actor
update. The update
(\href{https://github.com/verl-project/verl/blob/adc7eefa16dad75c5f7b878823d5a76eac90c7b3/verl/workers/engine_workers.py#L242-L310}{\texttt{train\_mini\_batch}})
splits the batch into mini-batches, passes over them \texttt{epochs} times, and takes one
optimizer step per mini-batch. Denote the parameters at these successive steps by
$\theta^{(0)}=\theta_{\mathrm{old}},\theta^{(1)},\theta^{(2)},\ldots$. At the next PPO
iteration, $\theta_{\mathrm{old}}$ is reset to the current parameters and the same pattern
repeats. Consequently, $w_t(\theta)=1$ holds only at the first optimizer step of each PPO
iteration; at every later step of the same iteration, $\theta\neq\theta_{\mathrm{old}}$.

In upstream VeRL the advantage is \emph{not} frozen at rollout time: the
\href{https://github.com/verl-project/verl/blob/adc7eefa16dad75c5f7b878823d5a76eac90c7b3/verl/trainer/distillation/losses.py#L373-L405}{per-token distillation loss}
is recomputed in every mini-batch from the current forward pass
(\texttt{model\_output["log\_probs"]}) and then detached. So at step $\theta^{(j)}$,
$\hat A_t=-\operatorname{sg}[\delta_t(y_t)]$ with $\delta_t$ evaluated at $\theta^{(j)}$.

\paragraph{Gradient of one term.}
Since $p_t^{\mathrm{old}}(y_t)$ does not depend on $\theta$, given
\eqref{eq:opd.log-derivative}, one has
\begin{align*}
    \nabla_\theta w_t(\theta)
    =
    \frac{\nabla_\theta p_t(y_t)}{p_t^{\mathrm{old}}(y_t)}
    =
    \frac{p_t(y_t)}{p_t^{\mathrm{old}}(y_t)}\,\nabla_\theta\log p_t(y_t)
    =
    w_t(\theta)\,\nabla_\theta\log p_t(y_t).
\end{align*}
Since $\hat A_t$ is constant in $\theta$ with value $-\delta_t(y_t)$, the gradient of
\eqref{eq:opd.ppo-clip} is
\begin{align*}
    -\hat A_t\,\nabla_\theta w_t(\theta)
    =
    w_t(\theta)\,\delta_t(y_t)\,\nabla_\theta\log p_t(y_t)
\end{align*}
when the clip is not binding, and zero when the clip (or the dual-clip bound) binds,
because the clipped branch is then constant in $\theta$.
\begin{itemize}
    \item \textbf{First optimizer step of each PPO iteration.} Here
    $\theta=\theta_{\mathrm{old}}$, so $w_t=1$ and neither clip binds, because
    $1\in[1-\epsilon_{\mathrm{low}},1+\epsilon_{\mathrm{high}}]$ and the dual-clip bound
    exceeds one. The gradient is $\delta_t(y_t)\nabla_\theta\log p_t(y_t)$, exactly the
    gradient of the sampled-token surrogate $\operatorname{sg}[\delta_t(y_t)]\log p_t(y_t)$ of
    Section~\ref{sec:opd.gradients}. Its expectation over $y_t\sim p_t$ is
    $\nabla_\theta\mathrm{kl}_t(x,y_{<t})$, i.e., term (A) of \eqref{eq:opd.exact-gradient}.
    \item \textbf{Later optimizer steps of the same iteration.} Now $w_t\neq 1$ in general,
    and $y_t$ was drawn from $p_t^{\mathrm{old}}$ rather than $p_t$. Where the clip does not
    bind, the ratio corrects for this exactly:
    \begin{align*}
        \E_{y_t\sim p_t^{\mathrm{old}}(\cdot\mid x,y_{<t})}\!\left[w_t\,\delta_t(y_t)\,\nabla_\theta\log p_t(y_t)\right]
        =
        \sum_{v\in\Vcal}p_t^{\mathrm{old}}(v)\,\frac{p_t(v)}{p_t^{\mathrm{old}}(v)}\,
        \delta_t(v)\,\nabla_\theta\log p_t(v)
        =
        \nabla_\theta\mathrm{kl}_t(x,y_{<t}),
    \end{align*}
    evaluated at the current $\theta$, where the last equality is \eqref{eq:opd.full-kl-grad}.
    This identity holds for a fixed prefix $(x,y_{<t})$; the expectation is over $y_t$ only.
    Two things still separate the update from term (A) of \eqref{eq:opd.exact-gradient} at
    the current $\theta$.
    \begin{enumerate}
        \item \emph{Prefixes.} Term (A) averages $\nabla_\theta\mathrm{kl}_t(x,y_{<t})$ over
        prefixes drawn from $\pi_\theta$, but these prefixes were drawn from
        $\pi_{\theta_{\mathrm{old}}}$. The ratio $w_t$ reweights only $y_t$, not $y_{<t}$.
        \item \emph{Clipping.} The identity uses the unclipped gradient. A token with
        $\hat A_t>0$ and $w_t>1+\epsilon_{\mathrm{high}}$, or with $\hat A_t<0$ and
        $w_t<1-\epsilon_{\mathrm{low}}$ (or beyond the dual-clip bound), contributes zero
        gradient, so the expected update differs from $\nabla_\theta\mathrm{kl}_t(x,y_{<t})$.
        This bias is intentional: it is PPO's trust region, which stops pushing a token's
        probability once it has moved far enough from the rollout policy.
    \end{enumerate}
\end{itemize}

\paragraph{Dictionary between PPO and OPD.}
\begin{itemize}
    \item \textbf{Reward.} The per-token reward is $-\delta_t(y_t)$: the teacher's
    log-probability minus the student's at the sampled token. There is no task reward.
    \item \textbf{Return and discount.} The advantage is the immediate reward only, i.e.,
    discount $\gamma=0$. This is exactly the step of dropping term (B) in
    \eqref{eq:opd.exact-gradient}. Using the reward-to-go
    $\hat A_t=-\sum_{t'\geq t}\operatorname{sg}[\delta_{t'}(y_{t'})]$ ($\gamma=1$) would add an
    unbiased estimate of term (B), since
    $\E[\delta_{t'}(y_{t'})\mid x,y_{<t'}]=\mathrm{kl}_{t'}(x,y_{<t'})$ and
    $\nabla_\theta\log p_t(y_t)$ is fixed given $y_{<t'}$ for $t<t'$.
    \item \textbf{Baseline and critic.} None; the advantage is not centered or normalized.
    \item \textbf{Ratio and clipping.} Inactive at the first optimizer step of each PPO
    iteration. At the later steps of the iteration, the ratio corrects the next-token
    sampling and the clip keeps $\theta$ close to $\theta_{\mathrm{old}}$, exactly as in PPO.
\end{itemize}
\begin{remark}[Is VeRL optimizing the OPD objective?]
    Ignore clipping. At later optimizer steps of a PPO iteration, the expected update is not
    term (A) of \eqref{eq:opd.exact-gradient} at the current $\theta$. This is not a bug, for
    three reasons.
    \begin{enumerate}
        \item \emph{It is the exact gradient of a frozen-state surrogate.} Define the objective
        with the prefix distribution frozen at the rollout policy,
        \begin{align*}
            \tilde J(\theta)
            \coloneqq
            \E_{(x,y)\sim\mu_{\theta_{\mathrm{old}}}}
            \left[
                \sum_{t=1}^{T(y)}\mathrm{kl}_t(x,y_{<t})
            \right],
        \end{align*}
        where only $\mathrm{kl}_t$, through $p_t$, depends on $\theta$. Its gradient is
        $\E_{\mu_{\theta_{\mathrm{old}}}}[\sum_t\nabla_\theta\mathrm{kl}_t(x,y_{<t})]$, which by the
        per-prefix identity above is the expectation of the unclipped update
        $w_t\,\delta_t(y_t)\,\nabla_\theta\log p_t(y_t)$ summed over $t$. So, with clipping
        ignored, VeRL runs unbiased stochastic gradient descent on $\tilde J$. In expectation,
        $\tilde J$ has the same gradient as the unclipped importance-weighted surrogate
        $\E_{\mu_{\theta_{\mathrm{old}}}}[\sum_t w_t(\theta)\hat A_t]$ used by TRPO; PPO optimizes a
        clipped version of that surrogate.
        \item \emph{The surrogate is accurate near $\theta_{\mathrm{old}}$.} At
        $\theta=\theta_{\mathrm{old}}$, $\nabla_\theta\tilde J$ equals term (A). For general
        $\theta$, the two differ only in the prefix distribution. Let
        $G\coloneqq\sup_{(x,y)}\|\sum_{t=1}^{T(y)}\nabla_\theta\mathrm{kl}_t(x,y_{<t})\|$, assumed
        finite (it is whenever $\Dcal_x$ has finite support, since responses are capped at
        $T_{\max}$). Since
        $\|\E_P[f]-\E_Q[f]\|\le 2\,\mathrm{TV}(P,Q)\sup\|f\|$, Pinsker's inequality, and the
        chain rule for the KL divergence give
        \begin{align*}
            \left\|
                \text{(A)}-\nabla_\theta\tilde J(\theta)
            \right\|
            \le
            2G\,\mathrm{TV}(\mu_\theta,\mu_{\theta_{\mathrm{old}}})
            \le
            G\sqrt{2\,D_{\mathrm{KL}}(\mu_{\theta_{\mathrm{old}}}\|\mu_\theta)},
            \\
            D_{\mathrm{KL}}(\mu_{\theta_{\mathrm{old}}}\|\mu_\theta)
            =
            \E_{\mu_{\theta_{\mathrm{old}}}}
            \left[
                \sum_{t=1}^{T(y)}D_{\mathrm{KL}}(p_t^{\mathrm{old}}\|p_t)
            \right].
        \end{align*}
        The gap is therefore controlled by how far the policy has moved since the rollout,
        which is what a trust region (PPO's clip, or TRPO's KL constraint) is designed to keep
        small. The prefix mismatch is thus not a separate flaw next to the trust region; it is
        the reason a trust region is needed.
        \item \emph{OPD already approximates more coarsely.} Dropping term (B), i.e., using
        $\gamma=0$, means that even the first optimizer step of each iteration does not follow
        the gradient of \eqref{eq:opd.token-kl}. Every standard OPD implementation makes this
        choice by design; the later optimizer steps add a second, local approximation on top
        of it.
    \end{enumerate}
    If exactness with respect to term (A) is required, use a single mini-batch and a single
    epoch per rollout. Every update is then an unbiased estimate of term (A), at the cost of
    sample efficiency.
\end{remark}

Rethinking-OPD, discussed next, uses the same machinery but replaces the single sampled
token $y_t$ by $k$ enumerated candidates.

\subsection{Rethink OPD modified VeRL}

\href{https://github.com/Thinking-Space/Rethinking-OPD/tree/ac26e38d6f1572eb027597b48a9f4e01f6915ef8}{Github Repo of Rethinking-OPD} implements a modified version of the student top-k OPD \eqref{eq:opd.top-k-objective}. It keeps the student
top-$k$ set $S_t$, but instead of backpropagating the renormalized divergence
$D_{\mathrm{KL}}(\bar p_t^{(S_t)}\|\bar q_t^{(S_t)})$, it runs the PPO update of
Section~\ref{sec:opd.as-ppo} with one advantage per \emph{candidate} $i\in S_t$ rather than
per sampled token. We first describe and analyze the method, then map it to the code.

\subsubsection{The Method}
\label{sec:opd.rethinking-method}

Throughout, $\delta_t(i)=\log p_t(i)-\log q_t(i)$ as in Section~\ref{sec:opd.gradients}.

\paragraph{Per-candidate advantage.}
Fix a prefix $(x,y_{<t})$. Every variant aims at the per-state gradient
$\nabla_\theta\mathrm{kl}_t(x,y_{<t})$ (term (A) of \eqref{eq:opd.exact-gradient}). Define the
per-token advantage function $\hat A_t(v)\coloneqq-\operatorname{sg}[\delta_t(v)]$ for
$v\in\Vcal$. By \eqref{eq:opd.full-kl-grad}, the target is an expectation under the student's
next-token distribution:
\begin{align}
    \nabla_\theta\mathrm{kl}_t(x,y_{<t})
    =
    \sum_{v\in\Vcal}p_t(v)\,\delta_t(v)\,\nabla_\theta\log p_t(v)
    =
    \E_{v\sim p_t}\!\left[-\hat A_t(v)\,\nabla_\theta\log p_t(v)\right].
    \label{eq:opd.kl-grad-expectation}
\end{align}
There are two natural ways to estimate this expectation.
\begin{itemize}
    \item \textbf{Sample one token (PG-OPD).} Draw $y_t\sim p_t$ and use
    $-\hat A_t(y_t)\nabla_\theta\log p_t(y_t)$. This is a one-sample Monte Carlo estimate of
    \eqref{eq:opd.kl-grad-expectation}: it is random, since it depends on $y_t$; it is
    unbiased; and it has sampling variance.
    \item \textbf{Enumerate a candidate set (Rethinking-OPD).} Replace $p_t$ by its
    renormalization $\bar p_t^{(S_t)}$ on $S_t$ and compute the expectation exactly:
    \begin{align*}
        \E_{i\sim\bar p_t^{(S_t)}}\!\left[-\hat A_t(i)\,\nabla_\theta\log p_t(i)\right]
        =
        -\sum_{i\in S_t}\bar p_t^{(S_t)}(i)\,\hat A_t(i)\,\nabla_\theta\log p_t(i).
    \end{align*}
    This estimate is deterministic given the prefix, so it has no variance over $y_t$. It
    equals \eqref{eq:opd.kl-grad-expectation} exactly when $S_t=\Vcal$ and is biased when
    $|S_t|<|\Vcal|$; the bias is characterized below.
\end{itemize}
The first estimate has the form ``advantage times the score $\nabla_\theta\log p_t(\cdot)$ of a
token''. The second estimate has the same form, once per candidate: the coefficient of
$\nabla_\theta\log p_t(i)$ is
\begin{align}
    A_{t,i}
    \coloneqq
    \bar p_t^{(S_t)}(i)\,\hat A_t(i)
    =
    -\operatorname{sg}\!\left[
        \bar p_t^{(S_t)}(i)\,\delta_t(i)
    \right],
    \qquad i\in S_t,
    \label{eq:opd.rethinking-cand-advantage}
\end{align}
so that the second estimate equals $-\sum_{i\in S_t}A_{t,i}\,\nabla_\theta\log p_t(i)$. That is,
$A_{t,i}$ is the PG-OPD advantage of token $i$, weighted by $\bar p_t^{(S_t)}(i)$. In PG-OPD this
weight is supplied implicitly by sampling; here the candidates are enumerated, so it must
appear explicitly. Because each candidate now looks like a sampled token with its own
advantage, the PPO loss of Section~\ref{sec:opd.as-ppo} can be reused unchanged, as described next.

\paragraph{The loss.}
The PPO objective \eqref{eq:opd.ppo-clip} is applied to each candidate and summed over
$S_t$. With the per-candidate ratio $w_{t,i}(\theta)\coloneqq p_t(i)/p_t^{\mathrm{old}}(i)$, the
per-state loss is
\begin{align}
    \sum_{i\in S_t}
    \max\!\Big\{
        -A_{t,i}\,w_{t,i}(\theta),\;
        -A_{t,i}\operatorname{clip}\!\big(w_{t,i}(\theta),1-\epsilon_{\mathrm{low}},1+\epsilon_{\mathrm{high}}\big)
    \Big\}.
    \label{eq:opd.rethinking-ppo}
\end{align}
The sampled token $y_t$ does not appear in \eqref{eq:opd.rethinking-ppo}; it only determines
which prefixes $(x,y_{<t})$ are visited. At the first optimizer step of each PPO iteration,
$w_{t,i}=1$, the clip is inactive, and $\nabla_\theta w_{t,i}=w_{t,i}\nabla_\theta\log p_t(i)=\nabla_\theta\log p_t(i)$,
exactly as in Section~\ref{sec:opd.as-ppo}. The gradient of \eqref{eq:opd.rethinking-ppo}
therefore equals that of the policy-gradient surrogate (``cand'' for candidate)
\begin{highlight}
    \begin{align}
        \ell_t^{\mathrm{cand}}
        \coloneqq
        -\sum_{i\in S_t}A_{t,i}\log p_t(i)
        =
        \sum_{i\in S_t}
        \operatorname{sg}\!\left[
            \bar p_t^{(S_t)}(i)\,\delta_t(i)
        \right]
        \log p_t(i).
        \label{eq:opd.rethinking-loss}
    \end{align}
\end{highlight}
Like the sampled-token surrogate, its value has no meaning; only its gradient matters.
With $S_t$ held fixed, that gradient is
\begin{align}
    g_t
    \coloneqq
    \nabla_\theta\ell_t^{\mathrm{cand}}
    =
    \sum_{i\in S_t}
    \bar p_t^{(S_t)}(i)\,\delta_t(i)\,\nabla_\theta\log p_t(i).
    \label{eq:opd.rethinking-grad}
\end{align}
At later optimizer steps of the same iteration, the per-candidate ratio and clip act
exactly as in Section~\ref{sec:opd.as-ppo}. The rest of this section analyzes $g_t$.

\paragraph{Relation to the full-vocabulary gradient.}
Compare $g_t$ with the exact per-state gradient \eqref{eq:opd.full-kl-grad}, derived in
step~3 of the sampled-token analysis in Section~\ref{sec:opd.gradients}. Sampled-token OPD estimates \eqref{eq:opd.full-kl-grad} with a single draw
$v=y_t\sim p_t$, so the sampling itself supplies the weight $p_t(v)$. Rethinking-OPD
enumerates the candidates instead of sampling them, so the weight must appear
explicitly: \eqref{eq:opd.rethinking-grad} is \eqref{eq:opd.full-kl-grad} truncated
to $S_t$ and renormalized, and it recovers \eqref{eq:opd.full-kl-grad} exactly when
$S_t=\Vcal$.

\paragraph{Comparison with student-top-$k$ OPD.}
Write $\bar\delta_t(i)\coloneqq\log\bar p_t^{(S_t)}(i)-\log\bar q_t^{(S_t)}(i)$ for the
renormalized log-ratio, so that
$D_{\mathrm{KL}}(\bar p_t^{(S_t)}\|\bar q_t^{(S_t)})=\sum_{i\in S_t}\bar p_t^{(S_t)}(i)\,\bar\delta_t(i)$.
The argument of step~3 in Section~\ref{sec:opd.gradients} applies verbatim to this pair of
distributions on $S_t$. It uses the product rule,
$\nabla_\theta\bar\delta_t(i)=\nabla_\theta\log\bar p_t^{(S_t)}(i)$,
\eqref{eq:opd.log-derivative}, and
\begin{align*}
    \sum_{i\in S_t}\bar p_t^{(S_t)}(i)\,\nabla_\theta\log\bar p_t^{(S_t)}(i)
    =
    \nabla_\theta\sum_{i\in S_t}\bar p_t^{(S_t)}(i)
    =
    \nabla_\theta 1
    =
    0.
\end{align*}
Placing the resulting top-$k$ gradient next to the other two per-state gradients,
\begin{align}
    \text{full vocabulary:}\quad
    \nabla_\theta\mathrm{kl}_t(x,y_{<t})
    &=
    \sum_{v\in\Vcal}p_t(v)\,\delta_t(v)\,\nabla_\theta\log p_t(v),
    \nonumber\\
    \text{student top-}k\text{:}\quad
    \nabla_\theta D_{\mathrm{KL}}\!\left(\bar p_t^{(S_t)}\,\middle\|\,\bar q_t^{(S_t)}\right)
    &=
    \sum_{i\in S_t}\bar p_t^{(S_t)}(i)\,\bar\delta_t(i)\,\nabla_\theta\log\bar p_t^{(S_t)}(i),
    \label{eq:opd.top-k-grad}\\
    \text{Rethinking-OPD:}\quad
    g_t
    &=
    \sum_{i\in S_t}\bar p_t^{(S_t)}(i)\,\delta_t(i)\,\nabla_\theta\log p_t(i).
    \nonumber
\end{align}
Top-$k$ OPD renormalizes everything: the weights, the log-ratio, and the score. It therefore
sees only the \emph{shape} of $p_t$ on $S_t$. Rethinking-OPD renormalizes only the weights and
keeps the unnormalized log-ratio and score of the full-vocabulary gradient.

To relate the last two, for $i\in S_t$, substitute
$\bar\delta_t(i)=\delta_t(i)-c_t$ with $c_t\coloneqq\log(P_{S_t}/Q_{S_t})$, where
$P_{S_t}=\sum_{i\in S_t}p_t(i)$ and $Q_{S_t}=\sum_{i\in S_t}q_t(i)$ are the student and teacher
masses on $S_t$ (as in the remark after \eqref{eq:opd.top-k-objective}), and
$\nabla_\theta\log\bar p_t^{(S_t)}(i)=\nabla_\theta\log p_t(i)-\nabla_\theta\log P_{S_t}$ into
\eqref{eq:opd.top-k-grad}. Using
$\sum_{i\in S_t}\bar p_t^{(S_t)}(i)\nabla_\theta\log p_t(i)=\sum_{i\in S_t}\nabla_\theta p_t(i)/P_{S_t}=\nabla_\theta\log P_{S_t}$
and $\sum_{i\in S_t}\bar p_t^{(S_t)}(i)=1$,
\begin{align*}
    \nabla_\theta D_{\mathrm{KL}}\!\left(\bar p_t^{(S_t)}\,\middle\|\,\bar q_t^{(S_t)}\right)
    &=
    \sum_{i\in S_t}\bar p_t^{(S_t)}(i)\,\big(\delta_t(i)-c_t\big)
    \big(\nabla_\theta\log p_t(i)-\nabla_\theta\log P_{S_t}\big)
    \\
    &=
    g_t
    -
    \E_{i\sim\bar p_t^{(S_t)}}\!\left[\delta_t(i)\right]\nabla_\theta\log P_{S_t}
    -
    c_t\,\nabla_\theta\log P_{S_t}
    +
    c_t\,\nabla_\theta\log P_{S_t}.
\end{align*}
The last two terms cancel, so
\begin{align}
    g_t
    =
    \underbrace{
    \nabla_\theta D_{\mathrm{KL}}\!\left(\bar p_t^{(S_t)}\,\middle\|\,\bar q_t^{(S_t)}\right)
    }_{\text{shape on } S_t}
    +
    \underbrace{
    \E_{i\sim\bar p_t^{(S_t)}}\!\left[\delta_t(i)\right]\nabla_\theta\log P_{S_t}
    }_{\text{mass on } S_t}.
    \label{eq:opd.rethinking-vs-topk}
\end{align}
The first term is exactly the gradient of student-top-$k$ OPD. The second term changes only the
student's total mass $P_{S_t}$ on $S_t$: under gradient descent it moves mass onto $S_t$ when
$\E_{\bar p}[\delta_t]<0$ (on average the teacher assigns more probability to the candidates
than the student) and off $S_t$ otherwise. This is the information that
\eqref{eq:opd.kl-decomposition} shows is discarded by renormalization. At $k=1$ the first term
vanishes and only the mass term remains, which is why Rethinking-OPD does not degenerate.

\paragraph{What $g_t$ descends.}
An update direction alone does not say what the method converges to. We therefore ask
whether $g_t$ is the gradient of some function at a fixed state. Recall the student
and teacher masses on $S_t$ from the remark after \eqref{eq:opd.top-k-objective},
$P_{S_t}=\sum_{i\in S_t}p_t(i)$ and $Q_{S_t}=\sum_{i\in S_t}q_t(i)$, and the
on-support part of the reverse KL in \eqref{eq:opd.kl-split},
\begin{align*}
    F_t
    \coloneqq
    \sum_{i\in S_t}p_t(i)\log\frac{p_t(i)}{q_t(i)}
    =
    \sum_{i\in S_t}p_t(i)\,\delta_t(i).
\end{align*}
By the product rule and \eqref{eq:opd.log-derivative}, with $\nabla_\theta q_t=0$,
\begin{align*}
    \nabla_\theta F_t
    =
    \sum_{i\in S_t}\big(\nabla_\theta p_t(i)\,\delta_t(i)+p_t(i)\,\nabla_\theta\log p_t(i)\big)
    =
    \sum_{i\in S_t}p_t(i)\,\delta_t(i)\,\nabla_\theta\log p_t(i)
    +
    \nabla_\theta P_{S_t}.
\end{align*}
Over the full vocabulary the last term would be $\nabla_\theta 1=0$, but on $S_t$ it is
not. Since $\bar p_t^{(S_t)}(i)=p_t(i)/P_{S_t}$, the first sum equals $P_{S_t}\,g_t$, so
\begin{align}
    g_t
    =
    \frac{1}{P_{S_t}}\,
    \nabla_\theta\Phi_t,
    \qquad
    \Phi_t
    \coloneqq
    F_t-P_{S_t}
    =
    \sum_{i\in S_t}
    \left[
        p_t(i)\log\frac{p_t(i)}{q_t(i)}-p_t(i)
    \right].
    \label{eq:opd.rethinking-potential}
\end{align}
Here $1/P_{S_t}$ is evaluated at the current parameters and not differentiated: it
enters through the detached weights and only rescales the step at each state. Hence,
at a fixed state and with $S_t$ fixed, Rethinking-OPD performs gradient descent on
$\Phi_t$, the on-support part of the reverse KL minus the student mass on $S_t$.

\paragraph{Minimizer of $\Phi_t$.}
Each summand of \eqref{eq:opd.rethinking-potential} depends on a single probability
$u=p_t(i)$ through $\varphi(u)=u\log(u/q_t(i))-u$, with $\varphi'(u)=\log(u/q_t(i))$ and
$\varphi''(u)=1/u>0$. So $\varphi$ is strictly convex and minimized at $u=q_t(i)$, and
$\Phi_t$ is minimized exactly when $p_t(i)=q_t(i)$ for all $i\in S_t$. This is
attainable, since $Q_{S_t}\le 1$ leaves mass $1-Q_{S_t}$ for the complement. Unlike
\eqref{eq:opd.top-k-objective}, the method therefore matches the teacher's
\emph{absolute} probabilities on $S_t$, not only their shape.\footnote{Adding the
constant $Q_{S_t}$, which does not change the gradient, gives
$\Phi_t+Q_{S_t}=\sum_{i\in S_t}[p_t(i)\log(p_t(i)/q_t(i))-p_t(i)+q_t(i)]\ge 0$, with
equality iff $p_t=q_t$ on $S_t$. This is the generalized KL divergence between the
unnormalized vectors $(p_t(i))_{i\in S_t}$ and $(q_t(i))_{i\in S_t}$.}
At $k=1$, $\Phi_t=\varphi(p_t(v_t^*))$ with $S_t=\{v_t^*\}$, minimized at
$p_t(v_t^*)=q_t(v_t^*)$, whereas \eqref{eq:opd.top-k-objective} is identically zero. At the
level of objectives, \eqref{eq:opd.kl-on-support} gives
$\Phi_t=P_{S_t}D_{\mathrm{KL}}(\bar p_t^{(S_t)}\|\bar q_t^{(S_t)})+\big[P_{S_t}\log(P_{S_t}/Q_{S_t})-P_{S_t}\big]$;
the bracket, as a function of $P_{S_t}$, is minimized at $P_{S_t}=Q_{S_t}$, which is the
objective-level counterpart of the mass term in \eqref{eq:opd.rethinking-vs-topk}.

\paragraph{Are the weights necessary?}
In sampled-token OPD the advantage $-\operatorname{sg}[\delta_t(y_t)]$ carries no probability
weight, so the weight $\bar p_t^{(S_t)}(i)$ in \eqref{eq:opd.rethinking-cand-advantage} may look
redundant. This is exactly the question raised in
\href{https://github.com/Thinking-Space/Rethinking-OPD/issues/13}{Rethinking-OPD issue \#13},
which proposes $A_{t,i}=-\operatorname{sg}[\delta_t(i)]$ (\icode{rm_scores = -kl_val}). The
weight is necessary: the candidates are enumerated rather than sampled, so nothing else
supplies the factor $p_t(v)$ of \eqref{eq:opd.full-kl-grad}. Without it the gradient becomes
\begin{align}
    \sum_{i\in S_t}\delta_t(i)\,\nabla_\theta\log p_t(i)
    =
    \nabla_\theta\,
    \frac{1}{2}\sum_{i\in S_t}\delta_t(i)^2,
\end{align}
an unweighted squared log-ratio on $S_t$. It is still minimized at $p_t=q_t$ on
$S_t$, but it is not a KL divergence. It also treats all $k$ candidates equally, so
tail candidates with small $p_t(i)$ and noisy log-ratios count as much as the head of
the distribution.


\subsubsection{Implementation Walkthrough}
\label{sec:opd.rethinking-impl}

The computation flow inside \icode{ppo/ray_trainer.py} is as follows. Each step computes one
of the quantities defined in Section~\ref{sec:opd.rethinking-method}.
\begin{enumerate}

    \item  
    \href{https://github.com/Thinking-Space/Rethinking-OPD/blob/ac26e38d6f1572eb027597b48a9f4e01f6915ef8/verl/verl/trainer/ppo/ray_trainer.py#L1051}
    {Generates the on-policy responses} via
    \newline
    \icode{old_log_probs, entropys, topk_ids, topk_log_probs = self.actor_rollout_wg.generate_sequences(gen_batch)}.

    Here, \icode{generate_sequences} produces the sampled student responses
    $y \sim \pi_\theta(\cdot\mid x)$ using the inference backend, e.g., vLLM or SGLang.

    \item Recomputes the \href{https://github.com/Thinking-Space/Rethinking-OPD/blob/ac26e38d6f1572eb027597b48a9f4e01f6915ef8/verl/verl/trainer/ppo/ray_trainer.py#L1111}{student log-probabilities} \icode{old_log_prob = self.actor_rollout_wg.compute_log_prob(batch)}.
    The training backend, e.g., FSDP, is used to compute the log-probabilities.
    The worker call eventually reaches
    \href{https://github.com/Thinking-Space/Rethinking-OPD/blob/ac26e38d6f1572eb027597b48a9f4e01f6915ef8/verl/verl/workers/actor/dp_actor.py#L648}
    {\texttt{DataParallelPPOActor.compute\_log\_prob}}, which invokes
    \newline
    \href{https://github.com/Thinking-Space/Rethinking-OPD/blob/ac26e38d6f1572eb027597b48a9f4e01f6915ef8/verl/verl/workers/actor/dp_actor.py#L86}{\texttt{DataParallelPPOActor.\_forward\_micro\_batch}}. When \texttt{top\_k > 0}, the actor
    computes the student top-$k$ token IDs and their log-probabilities:
    \begin{codeblock}[language=Python]
# assume input_ids is (B, T) and attention_mask is (B, T)
# attention_mask is 1 for real tokens and 0 for padding.
# NOTE: the input_ids contain (prefix, response, special tokens)
# no generation is performed here.

input_ids = input_ids.unsqueeze(-1) # (B, T) -> (B, T, 1)

# unpad_input removes all positions for which attention_mask == 0
# example: input_ids = [ [11, 12], [21, 0]], attention_mask = [ [1, 1], [1, 0]]
# then after unpad:
# input_ids_rmpad = [ [11], [12], [21]], indices = [0, 1, 3], cu_seqlens = [0, 2, 3]

input_ids_rmpad, indices, cu_seqlens, *_ = unpad_input(input_ids, attention_mask)
input_ids_rmpad = input_ids_rmpad.transpose(0, 1)  # (total_nnz, 1) -> (1, total_nnz)

# skip the details for 
# 1. unpadding the position_ids
# 2. handling the multi-modal inputs, and 
# 3. optionally Ulysses Sequence Parallelism

# compute the selected token ids by shifting the input_ids_rmpad to the right by
# 1 position. Example: input_ids_rmpad = [ [11], [12], [21]], then
# input_ids_rmpad_rolled = [ [12], [21], [11]]

input_ids_rmpad_rolled = torch.roll(input_ids_rmpad, shifts=-1, dims=1)
# (total_nnz, 1) -> (total_nnz,)
input_ids_rmpad_rolled = input_ids_rmpad_rolled.squeeze(0) 

# compute the logits; prevent model thinks we are generating
output = self.actor_module(
    input_ids=input_ids_rmpad,
    attention_mask=None,
    position_ids=position_ids_rmpad,
    ...
)

# get the logits for the full vocabulary
logits_rmpad = output.logits.squeeze(0)  # (total_nnz, vocab_size)
logits_rmpad.div_(temperature)

# compute the log-probabilities for the full vocabulary
log_probs_all = torch.log_softmax(logits_rmpad, dim=-1)  # (total_nnz, vocab_size)
# Gather log_probs for target tokens
log_probs = log_probs_all.gather(
    dim=-1, index=input_ids_rmpad_rolled.unsqueeze(-1)).squeeze(-1) # (total_nnz,)

# compute the top-k log-probabilities
_, topk_ids = torch.topk(logits_rmpad, k=top_k, dim=-1) # (total_nnz, top_k)

topk_log_probs = log_probs_all.gather(dim=-1, index=topk_ids) # (total_nnz, top_k)

full_topk_ids = pad_input(
    hidden_states=topk_ids,
    indices=indices,
    batch=batch_size,
    seqlen=seqlen,
)
full_topk_log_probs = pad_input(
    hidden_states=topk_log_probs,
    indices=indices,
    batch=batch_size,
    seqlen=seqlen,
)
# only return response part:
log_probs = full_log_probs.squeeze(-1)[:, -response_length - 1 : -1]  # (bsz, response_length)
topk_ids = full_topk_ids[:, -response_length - 1 : -1, :]
topk_log_probs = full_topk_log_probs[:, -response_length - 1 : -1, :]
    \end{codeblock}

    Therefore, for each student-visited state $s_t=(x,y_{<t})$,
    \[
        \texttt{student\_top\_k\_ids}
        =
        S_t
        =
        \operatorname{TopK}\!\left(
            \pi_\theta(\cdot\mid s_t)
        \right),
    \]
    and
    \[
        \texttt{student\_top\_k\_log\_probs}[i]
        =
        \log \pi_\theta(i\mid s_t) \equiv \log p_t(i)
        \qquad i\in S_t.
    \]

    % Importantly, the response tokens and the top-$k$ candidate tokens are
    % different objects. The response was sampled during
    % \texttt{generate\_sequences}; the top-$k$ set is computed by the actor
    % forward pass after generation.

    \item Then computes the
    \href{https://github.com/Thinking-Space/Rethinking-OPD/blob/ac26e38d6f1572eb027597b48a9f4e01f6915ef8/verl/verl/trainer/ppo/ray_trainer.py#L1135}
    {teacher log-probabilities} via the reward-model worker
    \newline
    \icode{teacher_data = self.rm_wg.compute_rm_score(batch)}. 
    
    Under the hood, it calls
    \href{https://github.com/Thinking-Space/Rethinking-OPD/blob/ac26e38d6f1572eb027597b48a9f4e01f6915ef8/verl/verl/workers/fsdp_workers.py#L2555}{\texttt{RewardModelWorker.compute\_rm\_score}}, which invokes 
    \newline
    \href{https://github.com/Thinking-Space/Rethinking-OPD/blob/ac26e38d6f1572eb027597b48a9f4e01f6915ef8/verl/verl/workers/fsdp_workers.py#L1925}{\texttt{RewardModelWorker.\_forward\_micro\_batch}}, which is very similar to \icode{DataParallelPPOActor._forward_micro_batch}. To highlight the core logic, we even use the most simple \href{https://github.com/Thinking-Space/Rethinking-OPD/blob/ac26e38d6f1572eb027597b48a9f4e01f6915ef8/verl/verl/workers/fsdp_workers.py#L2229}{branch}, i.e, not removing padding and not using Ulysses Sequence Parallelism.
    \begin{codeblock}[language=Python]
# the same sequence processing as in the actor forward pass
input_ids = micro_batch["input_ids"]        
output = self.reward_module(
    input_ids=input_ids,
    attention_mask=attention_mask,
    position_ids=position_ids,
    ...
)
rm_output_logits = output.logits
# rm_logits_resp: (bsz, response_length, vocab_size)
rm_logits_resp = rm_output_logits[:, -response_length - 1 : -1, :]  
rm_logits_resp = rm_logits_resp.div_(teacher_temperature)

# rm_log_probs: (bsz, response_length, vocab_size)
rm_log_probs = verl_F.logprobs_from_logits(
    logits=rm_logits_resp, labels=micro_batch["responses"])

# compute teacher top k log-probabilities

# rm_logits_resp: (bsz, response_length, vocab_size)
flat_logits = rm_logits_resp.reshape(-1, rm_logits_resp.size(-1)) # (N, vocab_size)
# student_top_k_ids: (bsz, seqlen, K_s)
flat_ids = student_top_k_ids.reshape(-1, student_top_k_ids.size(-1)) # (N, K_s)

flat_on_student_log_probs, flat_counts, flat_overlap_mask, 
    flat_teacher_top_k_ids, flat_teacher_top_k_log_probs, flat_teacher_in_student = 
    self._compute_teacher_top_k_log_probs(
    logits=flat_logits, student_ids=flat_ids, top_k=top_k, strategy=strategy)

# ... omit operations to get the tensors' shapess right
    \end{codeblock}
    At this point, the last piece is on the mechanics of \href{https://github.com/Thinking-Space/Rethinking-OPD/blob/ac26e38d6f1572eb027597b48a9f4e01f6915ef8/verl/verl/workers/actor/dp_actor.py#L1110}{\texttt{\_compute\_teacher\_top\_k\_log\_probs}}. It offers 5 options for selecting the top-k token sets:
    \begin{itemize}
        \item \texttt{only\_stu}: select top-k tokens from the student, then query the teacher for corresponding log-probs
        \item \texttt{only\_tch}: select top-k tokens from the teacher, then query the student for corresponding log-probs
        \item \texttt{intersection}: keep tokens appearing in both student and teacher top-k sets
        \item \texttt{union}: merge student and teacher top-k sets
        \item \texttt{union-intersection}: tokens in either Top-K but not both, i.e. symmetric difference
    \end{itemize}
    
    \begin{codeblock}[language=Python]
# we only show the code for the default strategy: only_stu

# logits: (n_samples, Vocab)
# student_ids: (n_samples, K_s)
# output: (n_samples, K_s)

# use chucked top_k to save the memory
for start in range(0, n_samples, chunk_size):
    end = min(start + chunk_size, n_samples)
    logits_chunk = logits[start:end] # (chunk, Vocab)
    student_ids_chunk = student_ids[start:end] # (chunk, K_s)
    
    # Top-K on Teacher
    # (chunk, K_t) and K_t = K_s = top_k
    t_logits, t_ids = torch.topk(logits_chunk, k=top_k, dim=-1) 
    t_logsumexp = torch.logsumexp(logits_chunk, dim=-1, keepdim=True)
    t_log_probs_top_k = t_logits - t_logsumexp

    # ... omitted for brevity ...
    
    chunk_log_probs = torch.gather(
        logits_chunk, dim=-1, index=student_ids_chunk) - t_logsumexp
    

    \end{codeblock}

    \item Once the student and teacher top-k log-probabilities are computed, the next step is to compute the reward. This is done by \href{https://github.com/Thinking-Space/Rethinking-OPD/blob/ac26e38d6f1572eb027597b48a9f4e01f6915ef8/verl/verl/trainer/ppo/ray_trainer.py#L1143}{\texttt{compute\_distillation\_reward}} 
    \newline
    \icode{distillation_output = self.actor_rollout_wg.compute_distillation_reward(batch)}, which calls
    \href{https://github.com/Thinking-Space/Rethinking-OPD/blob/ac26e38d6f1572eb027597b48a9f4e01f6915ef8/verl/verl/workers/actor/dp_actor.py#L451}{DataParallelPPOActor.compute\_distillation\_reward}. It also support the above 5 strategies. We use the default strategy \texttt{only\_stu} to demonstrate the computation.
    \begin{codeblock}[language=Python]
S_logp = student_top_k_log_probs # (n_samples, top_k)
T_logp = teacher_top_k_log_probs # (n_samples, top_k)
T_on_S = teacher_on_student_log_probs # (n_samples, top_k)

# normalize the student log-probabilities on the top-k support
log_probs = S_logp
norm_log_weights = log_probs - torch.logsumexp(log_probs, dim=-1, keepdim=True)
weights = torch.exp(norm_log_weights)
norm_weights = torch.nan_to_num(weights, nan=0.0, posinf=0.0, neginf=0.0)
kl_val = S_logp - T_on_S
rm_scores = -kl_val * norm_weights
    \end{codeblock}
    Here,
    \[
        \texttt{norm\_weights}_i
        =
        \frac{\pi_\theta(i\mid s_t)}
        {\sum_{j\in S_t}\pi_\theta(j\mid s_t)}
        \equiv \bar{p}_t^{(S_t)}(i)  \text{ for } i\in S_t
    \]
    and \texttt{kl\_val} is $\delta_t(i)$. The produced candidate-level reward is
    \begin{align}
        r_{t,i}
        =
        -\operatorname{sg}\!\left[
            \bar{p}_t^{(S_t)}(i)\,\delta_t(i)
        \right]
        \qquad \text{for } i\in S_t.
        \label{eq:opd.rethinking-reward}
    \end{align}
    i.e., the per-candidate advantage \eqref{eq:opd.rethinking-cand-advantage} before masking.
    The stop-gradient is automatic: \texttt{S\_logp} and \texttt{T\_on\_S} were
    computed in earlier forward passes and carry no gradient. Note that \texttt{rm\_scores}
    itself is never differentiated. Summing it over the candidates gives the value
    $-\sum_{i\in S_t}\bar p_t^{(S_t)}(i)\,\delta_t(i)$, but differentiating that sum through
    $\bar p_t^{(S_t)}$ would not give $g_t$; the gradient enters only through the PPO loss in
    the update step. Finally, \texttt{compute\_reward\_weights} has a \texttt{normalize=False}
    option (used by the \texttt{union-intersection} strategy), which uses the weights $p_t(i)$
    instead of $\bar p_t^{(S_t)}(i)$; the factor $1/P_{S_t}$ in
    \eqref{eq:opd.rethinking-potential} then disappears and the gradient is exactly
    $\nabla_\theta\Phi_t$.
    % At entry to this function, the provenance of its important inputs is
    % therefore
    % \begin{codeblock}[language=Python]
    %     # Produced by the student actor:
    %     S_ids  = data.batch["student_top_k_ids"]
    %     S_logp = data.batch["student_top_k_log_probs"]

    %     # Produced by the teacher/reward-model worker:
    %     T_on_S = data.batch["teacher_on_student_log_probs"]
    % \end{codeblock}

    % Thus,
    % \[
    %     \texttt{S\_logp}_i=\log\pi_\theta(i\mid s_t),
    %     \qquad
    %     \texttt{T\_on\_S}_i=\log\pi_T(i\mid s_t).
    % \]

    \item One the reward is computed, the next step is to compute the advantage. This is done by \href{https://github.com/Thinking-Space/Rethinking-OPD/blob/ac26e38d6f1572eb027597b48a9f4e01f6915ef8/verl/verl/trainer/ppo/ray_trainer.py#L1380}{\texttt{compute\_advantage}}, which ultimately calls \href{https://github.com/Thinking-Space/Rethinking-OPD/blob/ac26e38d6f1572eb027597b48a9f4e01f6915ef8/verl/verl/trainer/ppo/core_algos.py#L884}{\texttt{compute\_token\_reward\_direct\_advantage}}
\begin{codeblock}[language=Python]
advantage = compute_advantage(
    batch, # where you get the reward from
    adv_estimator= "token_reward_direct" # switch for OPD only.
    ... 
)
# in side: compute_token_reward_direct_advantage, rm_scores is reshaped to 
# (bsz, seqlen, top_k) and passed here as the token-level advantage 
advantage = rm_scores * response_mask
\end{codeblock}
    So the reward is used directly as the advantage, with no return, baseline, or
    normalization. Writing $m_t\in\{0,1\}$ for the response mask,
    \begin{align}
        A_{t,i}
        =
        m_t\, r_{t,i}
        =
        -m_t\operatorname{sg}\!\left[
            \bar{p}_t^{(S_t)}(i)\,\delta_t(i)
        \right]
        \qquad \text{for } i\in S_t.
        \label{eq:opd.rethinking-advantage}
    \end{align}
    On response positions this is exactly \eqref{eq:opd.rethinking-cand-advantage}.
    \item Lastly, the code reached to \href{https://github.com/Thinking-Space/Rethinking-OPD/blob/ac26e38d6f1572eb027597b48a9f4e01f6915ef8/verl/verl/trainer/ppo/ray_trainer.py#L2259}{\texttt{self.actor\_rollout\_wg.update\_actor}}, which is the entrypoint of the PPO update. If the FSDP backend is used, then the update done by \href{https://github.com/Thinking-Space/Rethinking-OPD/blob/ac26e38d6f1572eb027597b48a9f4e01f6915ef8/verl/verl/workers/fsdp_workers.py#L880}{self.actor.update\_policy}.
    With a 3D advantage, the actor recomputes $\log p_t(i)$ for the stored candidate
    ids, \icode{log_prob_for_loss = topk_log_probs}, and the
    \href{https://github.com/Thinking-Space/Rethinking-OPD/blob/ac26e38d6f1572eb027597b48a9f4e01f6915ef8/verl/verl/trainer/ppo/core_algos.py#L1118}{PPO loss}
    is applied per candidate and then summed over the candidates:
    \begin{codeblock}[language=Python]
# log_prob, old_log_prob, advantages: (bsz, response_length, top_k)
ratio = torch.exp(log_prob - old_log_prob)
pg_losses1 = -advantages * ratio
pg_losses2 = -advantages * torch.clamp(ratio, 1 - cliprange_low, 1 + cliprange_high)
pg_losses = torch.maximum(pg_losses1, pg_losses2)   # (+ dual-clip for A < 0)
pg_losses = torch.sum(pg_losses, dim=-1)            # sum over the top_k candidates
pg_loss = agg_loss(pg_losses, response_mask, loss_agg_mode)
    \end{codeblock}
    This is the candidate PPO loss \eqref{eq:opd.rethinking-ppo}, aggregated over response
    positions as specified by \texttt{loss\_agg\_mode}. The loop has the same structure as in
    Section~\ref{sec:opd.as-ppo}
    (\href{https://github.com/Thinking-Space/Rethinking-OPD/blob/ac26e38d6f1572eb027597b48a9f4e01f6915ef8/verl/verl/workers/actor/dp_actor.py#L787-L793}{mini-batches $\times$ \texttt{ppo\_epochs}},
    one optimizer step per mini-batch), with two differences. First, the advantage
    \eqref{eq:opd.rethinking-advantage} is computed once per PPO iteration from stored
    log-probabilities, so later optimizer steps of the same iteration use a stale
    $\delta_t(i)$ and $\bar p_t^{(S_t)}(i)$, unlike upstream VeRL. Second, when there is a
    single mini-batch and \texttt{ppo\_epochs=1} (the \texttt{on\_policy} flag), the code
    sets \icode{old_log_prob = log_prob_for_loss.detach()}, so $w_{t,i}\equiv 1$ by
    construction. At the first optimizer step of each PPO iteration (in particular, in the
    \texttt{on\_policy} case), the gradient is therefore $g_t$ of
    \eqref{eq:opd.rethinking-grad}, as established in Section~\ref{sec:opd.rethinking-method}.

\end{enumerate}

\subsubsection{Other unvisited reference}
\href{https://zoeyli.com/reinforcement%20learning/implementing-on-policy-distillation/}{Implementing On-Policy Distillation: Lessons from Building OPD in VeRL}




% There's one \href{https://github.com/verl-project/verl/blob/release/v0.9.0/verl/trainer/distillation/fsdp/losses.py}{implementation trick} for the top-k OPD: Compute selected log-probabilities without materializing an additional full [B,T,V] log-probability tensor.

% ## Discussion

% OPD 究竟在学什么: student 从自身当前的行为出发，通过不断访问真实状态，并在这些状态上对齐 teacher 的行为分布，从而吸收 teacher 在 RL 阶段形成的专精行为。

% OPD 理想的生效条件: 什么情况下，这种能力迁移会比较顺利？
% 1) 起点对齐 (解决的是“student 从哪里出发”的问题): 尽量让 student 在 OPD 开始时，接近各个 teacher 的 RL 起点分布。需要强调的是，这里的“接近”不意味着参数必须完全一致，也不意味着只要起点不同 OPD 就一定失败。更准确地说，它描述的是一种有利于能力迁移的初始化条件：起点越接近，通常越容易让 teacher 的后续行为变化被 student 有效吸收。
% 2) 数据对齐(解决的是“student 在哪里学习”的问题): OPD 阶段尽量直接复用 teacher RL 阶段使用过的 query，保持 query distribution 尽可能一致。原因很简单：teacher 的专精能力本身就是在这批 query 上通过 RL 训练和强化出来的。在相同的 query 上进行 OPD，student 能够直接进入 teacher 已经重点优化过的任务区域，更容易观察到 teacher 相对于自身当前策略所表现出来的专精优势。当然，从理论上说，并不能反过来推出“只有 teacher RL 训练过的 query 才能进行有效 OPD”。teacher 在 RL 阶段学到的能力可能具有一定的泛化性，因此，即使某个 query 没有直接出现在 teacher 的 RL 训练集中，只要它属于 teacher 专精能力能够发挥作用的任务分布，teacher 依然可能在这个 query 上表现出相对于 student 的能力优势。所以，这里需要区分理论上的可能性与工程上的最佳实践：理论上，OPD 不要求 query 必须是 teacher RL 训练集中的原始样本；但在实际工程中，为了最大程度复现 teacher 的 RL 能力、降低额外变量，我们通常直接复用 teacher RL 阶段的 query，并尽量保持数据分布一致。

% 当然，这两个条件并不是 OPD 成功的严格必要条件。它们更像是一个有助于提高迁移效率和稳定性的工程框架。真正的 OPD 成功与否，还取决于 teacher 是否确实提供了 student 尚未具备的有效能力、两者之间的能力差距是否处在可迁移范围内，以及训练过程本身是否稳定。

% 所以在 MOPD 的实现中，我们通常会：
% 各专项准备自己的 SFT 数据 → 统一混合 SFT → 所有 Teacher 从统一 SFT 启动 RL。
% 但是这个往往很难做到，所以现实中如果起点已经不一致，则通过轻量 SFT 尽量构造一个统一的近似起点。真正落到工程上，数据选择还有几个容易踩坑的地方。轻量 SFT 尽量复用 Teacher RL 起点对应的原始 SFT 数据，不随意换源，也不要用 Teacher RL 后的自蒸馏数据替代它。